\chapter{البيان الخامس: الاجتهاد والقياس وشروط المفتي}

\section{مقدمة في البيان الخامس (الاجتهاد)}
الحمد لله رب العالمين، والصلاة والسلام على نبينا محمد وعلى آله وصحبه أجمعين، أما بعد؛ فهذا هو الدرس الرابع من دروس شرح كتاب «الرسالة» للإمام أبي عبد الله الشافعي رحمه الله تعالى.

تحدثنا سابقاً عن وجوه البيان، ونصل اليوم إلى «البيان الخامس»، وهو بيان الاجتهاد؛ أي أن تجتهد للوصول إلى الحق عن طريق إعمال النص. وكعادة الإمام الشافعي في تأصيل القواعد، فإنه يكثر من ضرب الأمثلة لتتضح القاعدة، فذكر للاجتهاد أمثلة تطبيقية:

\section{أمثلة تطبيقية على الاجتهاد والقياس}

\subsection{الاجتهاد في تحديد جهة القبلة}
قال الله تبارك وتعالى: ﴿وَمِنْ حَيْثُ خَرَجْتَ فَوَلِّ وَجْهَكَ شَطْرَ الْمَسْجِدِ الْحَرَامِ﴾. الفرض على المسلم أينما كان أن يولي وجهه شطر المسجد الحرام. فإن كان معاينًا للكعبة (يراها أمامه)، فالفرض أن يتوجه لعينها يقيناً ولا مجال للاجتهاد. أما إن كان غائباً عنها، فالفرض أن يجتهد في التوجه إلى جهتها.

و«الشطر» في لغة العرب يأتي بمعنى النصف، ويأتي بمعنى الجهة، وهو المراد هنا. وقد استدل الشافعي بأشعار العرب على أن الشطر بمعنى الجهة (التلقاء والقصد)، كقول خفاف بن ندبة: «تُغني الرسالةُ شطرَ عمروٍ»، أي جهته. 
واجتهاد الغائب عن القبلة يكون بالدلائل والعلامات التي نصبها الله، قال تعالى: ﴿وَهُوَ الَّذِي جَعَلَ لَكُمُ النُّجُومَ لِتَهْتَدُوا بِهَا فِي ظُلُمَاتِ الْبَرِّ وَالْبَحْرِ﴾، وقال: ﴿وَعَلَامَاتٍ وَبِالنَّجْمِ هُمْ يَهْتَدُونَ﴾. فالمسلم يعمل عقله في هذه العلامات للوصول إلى جهة الكعبة، وهذا الإعمال للعلامات لمعرفة القبلة هو من باب الاجتهاد الصحيح.

\subsection{الاجتهاد في معرفة العدل}
ضرب الشافعي مثالاً آخر بالاجتهاد في تحديد العدول الذين تقبل شهادتهم وحكمهم. قال تعالى: ﴿وَأَشْهِدُوا ذَوَىْ عَدْلٍ مِّنكُمْ﴾. والعدل هو: من له مَلَكَةٌ تحمله على ملازمة التقوى والمروءة.
\begin{itemize}
    \item \textbf{التقوى:} اجتناب الشرك، والفسق، والبدعة، والإصرار على الصغيرة. وقال طلق بن حبيب: «التقوى أن تعمل بطاعة الله على نور من الله ترجو ثواب الله، وأن تترك معصية الله على نور من الله تخاف عقاب الله».
    \item \textbf{المروءة:} أن يلتزم الإنسان بما يليق بأمثاله قولا وفعلاً وسلوكاً، ويجتنب خوارم المروءة (كأن يأكل العالِم في الأسواق مثلاً).
\end{itemize}
فنحن نجتهد في تطبيق هذه الشروط على الأشخاص لنحكم بعدالتهم.

\subsection{الاجتهاد في جزاء الصيد (المثلية)}
المثال الثالث هو الاجتهاد في جزاء من قتل صيداً وهو محرم؛ قال تعالى: ﴿فَجَزَاءٌ مِّثْلُ مَا قَتَلَ مِنَ النَّعَمِ يَحْكُمُ بِهِ ذَوَا عَدْلٍ مِّنكُمْ﴾. بعد أن نجتهد في اختيار الحَكَمَين العدلين، يجتهد الحكمان في تحديد «المثل» من بهيمة الأنعام ليقابل الصيد البري المقتول. فإذا اصطاد حماراً وحشياً، حكما عليه ببقرة مققاربة له في الحجم، وإذا اصطاد ضبعاً حكما بكبش، وإذا اصطاد أرنباً حكما بعناق أو ما شابهه. ولا يعدل إلى «القيمة» المالية إلا إذا عُدم «المثل». وهذا الإلحاق بين الصيد والنعم هو عين «القياس الصحيح».

\section{مصادر التشريع وأنواع القياس عند الشافعي}
يقرر الشافعي أصلاً عظيماً: لا يجوز لأحد أبدًا أن يقول في شيء حلال أو حرام إلا من جهة العلم. وجهة العلم هي: الكتاب، أو السنة، أو الإجماع، أو القياس.

وبيّن أن القياس الصحيح يكون موافقاً للخبر، وأن الله ورسوله يحلان ويحرمان لأحد أمرين:
\begin{enumerate}
    \item \textbf{أمر تعبدي محض:} نص غير معقول المعنى (كعدد ركعات الصلاة، وتحريم الجمع بين المرأة وعمتها)، وهذا يوقف عنده ولا يقاس عليه.
    \item \textbf{معنى معقول (علّة):} أن يُحرم الشيء لعلة واضحة. فإذا وجدت هذه العلة في فرع لم ينص عليه، ألحقناه بالأصل. كالخمر حُرّمت لعلة الإسكار، فإذا وجد الإسكار في أي مشروب آخر (شعير، ذرة، حشيش) فهو حرام. وهذا يسمى \textbf{(قياس المعنى) أو (قياس العلة)}.
\end{enumerate}
وقد نجد فرعاً يشبه أصلين، فنلحقه بأولى الأشياء شبهاً به، وهذا يسمى \textbf{(قياس الشبه)}.

\section{بين الإجماع والاختلاف}
ينبه الشافعي إلى أدب عظيم لطالب العلم: ما أجمعت عليه الأمة لا يجوز الخلاف فيه (كحرمة الطواف بالقباب والأضرحة والذبح لها). أما ما اختلفت فيه الأمة اختلافاً سائغاً، فيُنظر في الدليل، ويُعمل بالأقرب للصواب. وقد قال الشافعي لمناظره يونس بن عبد الأعلى: «ألا يستقيم أن نكون إخواناً وإن لم نتفق في مسألة؟».
ولكن هذا في اختلاف الأفهام وتنوع الاستدلال، وليس في الاختلاف الذي يصادم النصوص الصريحة كمن يحلل ما حرم الله، أو يسوي بين المسلم والكافر، أو يدعو لمبادئ كفرية كالعلمانية والديمقراطية المستوردة المصادمة لحاكمية الشريعة بدعوى الخلاف، فهذا اختلاف تضاد لا يقبل.

\section{شروط المفتي والمجتهد}
بيّن الشافعي أن المجتهد يجب أن تتوفر فيه علوم أساسية، أهمها:

\subsection{العلم بلغة العرب}
فالقرآن نزل بلسان عربي مبين، ولغة العرب هي لغة التشريع. والصحابة كانوا يفهمون دلالات الألفاظ بالسليقة. فعلى المفتي أن يعلم العام، والخاص، والمطلق، والمقيد، والعام الذي يراد به الخصوص. والاعتزاز بالعربية هو جزء من الاعتزاز بالهوية الإسلامية. وقد زعم بعض المقلدين وجود كلمات أعجمية في القرآن، ورد الشافعي ذلك بقوة، مقرراً أن كل ما في القرآن فهو عربي خالص، وما وافق لغات أخرى فهو إما من توارد اللغات، أو كلمات استعملتها العرب وعرّبتها قبل نزول القرآن فصارت من لسانها.

\subsection{العلم بالناسخ والمنسوخ ودلالات الأوامر}
يجب أن يعلم المفتي الناسخ والمنسوخ (ومنه ما نُسخ حكماً وبقي تلاوة، أو نُسخ تلاوة وبقي حكماً كآية الرجم). كما يجب أن يعلم دلالة الأوامر والنواهي؛ فمن الأوامر ما يكون للوجوب، ومنها ما يكون للندب، ومنها ما يكون للاباحة.

\section{التحذير من القول بلا علم ومن التقليد الأعمى}
ختم الشافعي هذا الباب بكلمة تكتب بماء الذهب: \textbf{«فالواجب على العالمين ألّا يقولوا إلا من حيث علموا. وقد تكلم في العلم من لو أمسك عن بعض ما تكلم فيه منه، لكان الإمساك أولى به، وأقرب من السلامة له إن شاء الله»}.
فلا يجوز للعالم أن يتكلم إلا بعلم راسخ من الكتاب والسنة. كما حذر من التقليد الأعمى، وعبر عنه بقوله: «وبالتقليد أغفل من أغفل منهم، والله يغفر لنا ولهم». فالمقلد الأعمى يتخبط في الظلمات، والواجب على طالب العلم إذا بلغه قول أن يسأل عن دليله الشرعي.